% ============================================================================
%  SECCIÓN 3.28: HISTORIAS DE USUARIO
% ============================================================================

\section{Historias de usuario}

Las historias de usuario permiten capturar los requisitos desde la perspectiva
del usuario final, describiendo \textit{qué} necesita y \textit{por qué}, sin
especificar \textit{cómo} se implementará \citep{beck2001, cohn2004}. A
continuación se presentan las historias de usuario definidas para el MVP de
la aplicación \nombreApp{}, organizadas por épico.

\subsection{Épico: Autenticación}

\begin{table}[H]
  \centering
  \caption{Historias de usuario -- Autenticación}
  \contenidoConFuente{%
    \begin{tablaAPA}
      \begin{tabular}{@{}p{1.2cm}m{7.0cm}m{2.0cm}m{2.5cm}m{2.5cm}@{}}
        \toprule
        \textbf{ID} & \textbf{Historia} & \textbf{Prioridad} &
        \textbf{Estado} & \textbf{Sprint} \\
        \midrule
        HU-01 & Como visitante, quiero crear una cuenta con mi correo y
        contraseña para poder acceder a la aplicación. & Alta & Completada &
        Sprint 1 \\
        \addlinespace
        HU-02 & Como usuario, quiero iniciar sesión con mis credenciales
        para acceder a mi perfil y datos. & Alta & Completada & Sprint 1 \\
        \addlinespace
        HU-03 & Como usuario, quiero cerrar sesión de forma segura para
        proteger mis datospersonales. & Alta & Completada & Sprint 1 \\
        \addlinespace
        HU-04 & Como usuario, quiero que la app recuerde mi sesión para no
        tener que iniciar sesión cada vez que la abra. & Media & Completada &
        Sprint 1 \\
        \bottomrule
      \end{tabular}
    \end{tablaAPA}%
  }{Elaboración propia.}
\end{table}

\subsection{Épico: Exploración de escenarios}

\begin{table}[H]
  \centering
  \caption{Historias de usuario -- Exploración}
  \contenidoConFuente{%
    \begin{tablaAPA}
      \begin{tabular}{@{}p{1.2cm}m{7.0cm}m{2.0cm}m{2.5cm}m{2.5cm}@{}}
        \toprule
        \textbf{ID} & \textbf{Historia} & \textbf{Prioridad} &
        \textbf{Estado} & \textbf{Sprint} \\
        \midrule
        HU-05 & Como usuario, quiero ver el listado de escenarios
        deportivos para conocer la oferta disponible. & Alta & Completada &
        Sprint 2 \\
        \addlinespace
        HU-06 & Como usuario, quiero buscar escenarios por nombre para
        encontrar rápidamente uno específico. & Alta & Completada & Sprint 2
        \\
        \addlinespace
        HU-07 & Como usuario, quiero filtrar escenarios por deporte y tipo
        para reducir los resultados a mi interés. & Alta & Completada &
        Sprint 2 \\
        \addlinespace
        HU-08 & Como usuario, quiero ver la ubicación de los escenarios en
        un mapa para planificar mi visita. & Alta & Completada & Sprint 2 \\
        \addlinespace
        HU-09 & Como usuario, quiero ver la ubicación actual en el mapa
        para orientarme respecto a los escenarios cercanos. & Media &
        Completada & Sprint 2 \\
        \addlinespace
        HU-10 & Como usuario, quiero ver el detalle completo de un
        escenario (disciplinas, horarios, servicios) para decidir si lo
        visito. & Alta & Completada & Sprint 2 \\
        \addlinespace
        HU-11 & Como usuario, quiero ver los eventos programados en un
        escenario para planificar mi asistencia. & Media & Completada &
        Sprint 2 \\
        \bottomrule
      \end{tabular}
    \end{tablaAPA}%
  }{Elaboración propia.}
\end{table}

\subsection{Épico: Favoritos}

\begin{table}[H]
  \centering
  \caption{Historias de usuario -- Favoritos}
  \contenidoConFuente{%
    \begin{tablaAPA}
      \begin{tabular}{@{}p{1.2cm}m{7.0cm}m{2.0cm}m{2.5cm}m{2.5cm}@{}}
        \toprule
        \textbf{ID} & \textbf{Historia} & \textbf{Prioridad} &
        \textbf{Estado} & \textbf{Sprint} \\
        \midrule
        HU-12 & Como usuario, quiero guardar un escenario como favorito
        para acceder rápidamente después. & Alta & Completada & Sprint 2 \\
        \addlinespace
        HU-13 & Como usuario, quiero quitar un escenario de mis favoritos
        si ya no me interesa. & Alta & Completada & Sprint 2 \\
        \addlinespace
        HU-14 & Como usuario, quiero ver mi lista de favoritos para
        consultar los escenarios que guardé. & Alta & Completada & Sprint 2
        \\
        \bottomrule
      \end{tabular}
    \end{tablaAPA}%
  }{Elaboración propia.}
\end{table}

\subsection{Épico: Gestión de contenido}

\begin{table}[H]
  \centering
  \caption{Historias de usuario -- Gestión}
  \contenidoConFuente{%
    \begin{tablaAPA}
      \begin{tabular}{@{}p{1.2cm}m{7.0cm}m{2.0cm}m{2.5cm}m{2.5cm}@{}}
        \toprule
        \textbf{ID} & \textbf{Historia} & \textbf{Prioridad} &
        \textbf{Estado} & \textbf{Sprint} \\
        \midrule
        HU-15 & Como gestor, quiero registrar la información de mi
        escenario para darle visibilidad. & Media & Pendiente & Sprint 3 \\
        \addlinespace
        HU-16 & Como gestor, quiero actualizar los datos de mi escenario
        para mantener la información al día. & Media & Pendiente & Sprint 3
        \\
        \addlinespace
        HU-17 & Como administrador, quiero gestionar los usuarios del
        sistema para validar la información publicada. & Baja & Pendiente &
        Futuro \\
        \bottomrule
      \end{tabular}
    \end{tablaAPA}%
  }{Elaboración propia.}
\end{table}

\subsection{Épico: Interfaz y experiencia}

\begin{table}[H]
  \centering
  \caption{Historias de usuario -- Interfaz}
  \contenidoConFuente{%
    \begin{tablaAPA}
      \begin{tabular}{@{}p{1.2cm}m{7.0cm}m{2.0cm}m{2.5cm}m{2.5cm}@{}}
        \toprule
        \textbf{ID} & \textbf{Historia} & \textbf{Prioridad} &
        \textbf{Estado} & \textbf{Sprint} \\
        \midrule
        HU-18 & Como usuario, quiero ver una pantalla de bienvenida con el
        logo para identificar la app. & Media & Completada & Sprint 1 \\
        \addlinespace
        HU-19 & Como usuario, quiero que la app tenga un tema visual
        consistente para una experiencia agradable. & Media & Completada &
        Sprint 3 \\
        \addlinespace
        HU-20 & Como usuario, quiero que la app muestre indicadores de
        carga mientras obtiene datos para saber que está funcionando. &
        Media & Completada & Sprint 2 \\
        \addlinespace
        HU-21 & Como usuario, quiero ver mensajes amigables cuando no hay
        datos disponibles para entender la situación. & Baja & Completada &
        Sprint 2 \\
        \addlinespace
        HU-22 & Como usuario, quiero refrescar el contenido con un gesto
        de deslizar para obtener datos actualizados. & Baja & Completada &
        Sprint 2 \\
        \bottomrule
      \end{tabular}
    \end{tablaAPA}%
  }{Elaboración propia.}
\end{table}

\subsection{Mapa de historias de usuario}

La siguiente figura presenta la relación entre las historias de usuario y los
módulos de la aplicación:

\begin{figure}[H]
  \centering
  \caption{Mapa de historias de usuario por módulo}
  \label{fig:mapa-historias}
  \contenidoConFuente{%
  \begin{tikzpicture}[
      modulo/.style={rectangle, draw=black!60, rounded corners=2pt,
                     align=center, fill=blue!8, font=\scriptsize,
                     text width=3.2cm, minimum height=1.0cm},
      historia/.style={rectangle, draw=black!40, rounded corners=1pt,
                       align=left, fill=green!5, font=\tiny,
                       text width=2.8cm, minimum height=0.5cm},
      flecha/.style={-{Stealth[length=1.5mm]}, thick, draw=black!40}
    ]
    \node[modulo] (auth) at (0,3) {Autenticación\\(HU-01 a HU-04)};
    \node[modulo] (expl) at (-4.5,1) {Exploración\\(HU-05 a HU-11)};
    \node[modulo] (fav) at (4.5,1) {Favoritos\\(HU-12 a HU-14)};
    \node[modulo] (gest) at (-4.5,-1.5) {Gestión\\(HU-15 a HU-17)};
    \node[modulo] (ui) at (4.5,-1.5) {Interfaz\\(HU-18 a HU-22)};

    \draw[flecha] (auth) -- (-1.5,2);
    \draw[flecha] (auth) -- (1.5,2);
    \draw[flecha] (expl) -- (-1.5,0.5);
    \draw[flecha] (fav) -- (1.5,0.5);
  \end{tikzpicture}%
  }{Elaboración propia.}
\end{figure}

\subsection{Criterios de aceptación (ejemplo)}

A continuación se presentan los criterios de aceptación para las historias de
usuario más representativas:

\subsubsection{HU-05: Ver listado de escenarios}

\begin{enumerate}
  \item Al abrir la pantalla de catálogo, se muestra el listado completo de
    escenarios deportivos.
  \item Cada elemento del listado muestra: nombre, dirección, deportes
    disponibles e imagen principal.
  \item El listado se carga en menos de 3 segundos.
  \item Si no hay escenarios disponibles, se muestra un mensaje informativo.
  \item Se puede deslizar hacia abajo para refrescar los datos.
\end{enumerate}

\subsubsection{HU-06: Buscar escenarios por nombre}

\begin{enumerate}
  \item Existe un campo de búsqueda en la parte superior de la pantalla de
    catálogo.
  \item Al escribir al menos 2 caracteres, se filtran los escenarios cuyo
    nombre contiene el texto ingresado.
  \item La búsqueda es en tiempo real (se actualiza con cada carácter
    ingresado).
  \item Si no hay resultados, se muestra un mensaje indicando que no se
    encontraron escenarios.
\end{enumerate}

\subsubsection{HU-08: Ver ubicación en el mapa}

\begin{enumerate}
  \item Al abrir la pantalla de mapa, se muestran los marcadores de todos
    los escenarios.
  \item Al tocar un marcador, se muestra un popup con el nombre y dirección
    del escenario.
  \item El mapa muestra la ubicación actual del usuario con un marcador
    especial.
  \item El mapa permite hacer zoom y desplazarse libremente.
\end{enumerate}

\subsubsection{HU-12: Guardar escenario como favorito}

\begin{enumerate}
  \item En la pantalla de detalle, existe un botón de ``favorito'' (ícono
    de estrella o corazón).
  \item Al tocar el botón, el escenario se agrega a la lista de favoritos
    del usuario.
  \item El botón cambia de apariencia para indicar que el escenario ya es
    favorito.
  \item El favorito se persiste en la base de datos y sobrevive al cierre
    de la aplicación.
\end{enumerate}
